% Generated draft: requires mathematical and source review.
\documentclass[11pt]{article}
\usepackage[margin=25mm]{geometry}
\usepackage{fontspec}
\setmainfont{Latin Modern Roman}
\usepackage{amsmath,amssymb,mathtools,unicode-math}
\setmathfont{Latin Modern Math}
\usepackage{enumitem,needspace,xurl,hyperref}
\setlength{\parindent}{0pt}
\setlength{\parskip}{5pt}
\setlength{\emergencystretch}{3em}
\newcommand{\N}{\mathbb{N}}
\newcommand{\Z}{\mathbb{Z}}
\newcommand{\Q}{\mathbb{Q}}
\newcommand{\R}{\mathbb{R}}
\newcommand{\C}{\mathbb{C}}
\newcommand{\F}{\mathbb{F}}
\newcommand{\A}{\mathbb{A}}
\newcommand{\PP}{\mathbb{P}}
\newcommand{\E}{\mathbb{E}}
\newcommand{\eps}{\varepsilon}
\newcommand{\dd}{\,\mathrm{d}}
\begin{document}
\section*{UnsolvedMath ID: 30005058}
\textbf{Source problem code:} catalogue entry
\subsection{Definitions and mathematical statement}
For interpreting the displayed comparison, let \(k,m\) be positive integers, let \(X_i\) be random vectors taking values in \(\R^{n_i}\), and write \(X=(X_1,\ldots,X_k)\). Under the conventional coupling interpretation, \(\Pi(X_1,\ldots,X_k)\) denotes the collection of joint probability laws whose respective marginals are the laws of \(X_1,\ldots,X_k\). The supplied excerpt does not explicitly define \(\Pi\), so this interpretation requires confirmation. For a random vector \(Y\) with Lebesgue density \(f_Y\), its differential entropy, when well-defined and finite, is
\[
h(Y)=-\int f_Y(y)\log f_Y(y)\,\dd y.
\]
Let \(d_j\geq 0\) and let \(B_j:\R^{n_1+\cdots+n_k}\to\R^{r_j}\) be linear maps, for \(1\leq j\leq m\). These domains express the dimensional compatibility mentioned in the excerpt; the remaining hypotheses of the cited principle are not supplied.

The source fragment describes a Gaussian comparison in which Gaussian random vectors \(Z_1,\ldots,Z_k\) satisfy
\[
\dim Z_i=\dim X_i,\qquad h(Z_i)=h(X_i)\quad(1\leq i\leq k),
\]
and displays the schematic inequality
\begin{equation}
\max_{X\in\Pi(X_1,\ldots,X_k)}\sum_{j=1}^{m}d_jh(B_jX)
\;\geq\;
\max_{Z\in\Pi(Z_1,\ldots,Z_k)}\sum_{j=1}^{m}d_jh(B_jZ).
\tag{1}
\end{equation}
Here membership in \(\Pi\) is interpreted as a condition on the law of the concatenated vector. Gaussian marginal laws do not by themselves require a jointly Gaussian coupling. The source fragment does not specify whether the right-hand optimization has an additional joint Gaussian restriction. The maxima in (1) are retained as supplied; their existence, entropy finiteness, admissible couplings, and any nondegeneracy or regularity hypotheses are unresolved. Consequently, (1) is recorded as source context, not asserted here as a theorem for every choice of the displayed data.

\textbf{Open research target, as supplied.} Find a common Gaussian-comparison framework that encompasses the entropic and geometric inequalities unified by (1) and also proves the additional inequalities referred to in the report as not obviously implied by (1). The supplied record does not identify those additional inequalities or give a precise criterion for what constitutes a common framework. Thus the target remains a research-program question rather than a fully specified quantified conjecture. In particular, ``not obviously implied'' does not assert that implication from (1) is impossible.

The supplied curation labels the problem partially solved and attributes a broad unifying Gaussian-comparison framework to Aras and Courtade, while reporting incomplete coverage of the additional inequalities. These are supplied literature assessments, not independently verified mathematical conclusions.
\subsection{Short English statement}
Seek a common Gaussian-comparison framework covering both the inequalities unified by the report's comparison principle and the additional inequalities it describes as not obviously implied. The omitted hypotheses of that principle and the unspecified additional inequalities prevent a fully precise conjecture from being extracted from this record.
\subsection{Sources}
\begin{itemize}
\item Supplied local registry record: \texttt{data/unsolvedmath/problems.json}, permanent UnsolvedMath ID \(30005058\), exact source problem code \texttt{catalogue entry}. No individual proposer is identified; the supplied proposed year is 2022.
\item Supplied report citation: \emph{New Mathematical Techniques in Information Theory} (2023), Oberwolfach Reports, pp.\ 683--707, DOI \texttt{10.4171/owr/2022/14}. Supplied literature-source label: ``Oberwolfach Report 14/2022.'' \url{https://doi.org/10.4171/owr/2022/14}. The report is supplied as the authoritative source; its full text has not been independently consulted for this draft. The citation's 2023 date and the record's 2022 report designation are retained without resolving their relationship.
\item Supplied literature reference and attribution: Aras--Courtade, ``entropy inequalities and Gaussian comparisons.'' \url{https://arxiv.org/abs/2206.14182}. This is the supplied reference label; a fuller bibliographic citation is not provided.
\item Literature assessment date: \texttt{2026-08-22}. This is supplied dated curation, not independently verified. It labels the problem partially solved, credits Aras and Courtade with a broad framework unifying many entropy and geometric inequalities, and reports that the broader coverage requested by the report remains incomplete. The supplied background explicitly requests independent verification of time-sensitive status and cited claims.
\item The registry labels the statement \texttt{corrected\_verified} and reports a check against the official Oberwolfach report. This is supplied verification metadata, not verification performed for this draft.
\end{itemize}
\end{document}
